\begin{notechapter}
\chapter{Bigram Language Models and Softmax Regression}\label{ch:bigram-models}
\chaptersubtitle{Counting, the softmax on a one-hot input, and why the two agree}

\begin{chapterintro}
A language model assigns a probability to every sequence of words. The simplest
non-trivial one, the \emph{bigram} model, predicts each word from the previous word only.
It can be estimated by counting, or trained as a multinomial logistic regression (softmax
regression) whose input is the one-hot code of the previous word. This chapter shows that
the two approaches are the same model, derives the gradient once and carefully, and
uses the comparison to explain what the one-hot target of the cross-entropy loss means,
why unseen bigrams are a problem, and what regularization does about it. You will learn
\begin{itemize}[nosep]
  \item the bigram model as an autoregressive model with a one-word memory, and its
        maximum-likelihood estimate by counting;
  \item softmax regression on a one-hot input: the product $W^\top x$ is a row lookup;
  \item the gradient of the cross-entropy, from the logits to the matrix $W$;
  \item why training on one-hot targets is not an approximation of the ``true''
        distribution of the next word;
  \item why the trained model converges to the counts, what happens to unseen bigrams, and
        how perplexity measures a language model.
\end{itemize}
\prerequisites{the softmax and the cross-entropy (\cref{ch:sigmoid-softmax,ch:cross-entropy});
backpropagation through an affine map (\cref{backpropagation:sec:affine}); autoregressive
models (\cref{sequence-models:sec:autoregressive}).}
\end{chapterintro}

\section{Language models and the bigram assumption}\label{bigram-models:sec:bigram}

Let the vocabulary be $\{1,\dots,V\}$, with two special symbols: $\langle s\rangle$, which starts every sentence, and
$\langle/s\rangle$, which ends it. A sentence is a sequence $w_0 = \langle s\rangle, w_1,\dots,w_T = \langle/s\rangle$. By the chain rule of
probability,
\[
  p(w_1,\dots,w_T) = \prod_{t=1}^Tp(w_t\mid w_0,\dots,w_{t-1}),
\]
so a language model is an autoregressive model (\cref{sequence-models:sec:autoregressive}): a
model of the next word given all the previous ones. The end symbol makes the probabilities
of sentences of all lengths sum to $1$. A recurrent network can condition on the whole prefix.
The \emph{bigram} model makes the Markov assumption that only the last word matters:
\[
  p(w_t\mid w_0,\dots,w_{t-1}) = p(w_t\mid w_{t-1}) .
\]
It is described by a $V\times V$ table $\theta_{ij} = p(\text{next} = j\mid\text{previous} = i)$, each row a probability distribution.

\paragraph{Maximum likelihood by counting.} Concatenate the training sentences into $N$
bigrams $(i_n, j_n)$, and let $C_{ij}$ be the number of occurrences of the bigram $(i, j)$ and $C_i = \sum_jC_{ij}$ the
number of bigrams that start with $i$.

\begin{proposition}[Count estimate]\label{bigram-models:prop:counts}
The log-likelihood $\sum_n\log\theta_{i_nj_n} = \sum_{i,j}C_{ij}\log\theta_{ij}$ is maximized, over all tables whose rows are
probability distributions, by $\hat\theta_{ij} = C_{ij}/C_i$ for every $i$ with $C_i > 0$. Rows with $C_i = 0$ are not determined.
\end{proposition}

\begin{proof}
The log-likelihood is a sum of one term per row, $\sum_jC_{ij}\log\theta_{ij}$, and the rows are constrained
separately, so each row is maximized separately. For $C_i > 0$, $\sum_jC_{ij}\log\theta_{ij} = C_i\sum_j\hat\theta_{ij}\log\theta_{ij}$, and by Gibbs'
inequality ($\mathrm{KL}(\hat\theta_{i\cdot}\Vert\theta_{i\cdot})\geq0$, \cref{ch:cross-entropy}) $\sum_j\hat\theta_{ij}\log\theta_{ij}\leq\sum_j\hat\theta_{ij}\log\hat\theta_{ij}$, with equality iff $\theta_{i\cdot} = \hat\theta_{i\cdot}$.
\end{proof}

\begin{example}[A four-sentence corpus]\label{bigram-models:ex:corpus}
Take the sentences \emph{the cat sat}, \emph{the dog sat}, \emph{the cat ran}, \emph{a cat sat}, with the vocabulary
$\langle s\rangle$, the, a, cat, dog, sat, ran, $\langle/s\rangle$ ($V = 8$). Each sentence gives four bigrams, $N = 16$. The non-zero
estimates are
\begin{gather*}
  \hat p(\text{the}\mid\langle s\rangle) = \tfrac34,\quad \hat p(\text{a}\mid\langle s\rangle) = \tfrac14,\qquad
  \hat p(\text{cat}\mid\text{the}) = \tfrac23,\quad \hat p(\text{dog}\mid\text{the}) = \tfrac13,\\
  \hat p(\text{sat}\mid\text{cat}) = \tfrac23,\quad \hat p(\text{ran}\mid\text{cat}) = \tfrac13,
\end{gather*}
and $\hat p(\text{cat}\mid\text{a}) = \hat p(\text{sat}\mid\text{dog}) = \hat p(\langle/s\rangle\mid\text{sat}) = \hat p(\langle/s\rangle\mid\text{ran}) = 1$. The row of $\langle/s\rangle$ is empty: nothing follows the end of a
sentence.
\end{example}

\section{Softmax regression on a one-hot input}\label{bigram-models:sec:softmax}

The same model can be written as a one-layer network. Encode the previous word $i$ by the
one-hot vector $x = e_i\in\R^V$ and compute
\[
  z = W^\top x + b,\qquad p = \softmax(z),\qquad p_j = \frac{e^{z_j}}{\sum_{k=1}^Ve^{z_k}} = p_\theta(\text{next} = j\mid\text{previous} = i),
\]
with $W\in\R^{V\times V}$ and $b\in\R^V$; $z$ is the vector of \emph{logits}. Here $W$ multiplies $x$ transposed, as in the
original notes and in \texttt{x @ W} with row vectors; the layer of \cref{ch:backpropagation} would
be written with $W^\top$ in place of $W$.

\paragraph{A row lookup.} For a one-hot input, the product selects one row:
\[
  W^\top e_i = \begin{bmatrix}w_{11} & w_{21} & \cdots & w_{V1}\\ \vdots & \vdots & & \vdots\\ w_{1V} & w_{2V} & \cdots & w_{VV}\end{bmatrix}
  \begin{bmatrix}0\\ \vdots\\ 1\\ \vdots\\ 0\end{bmatrix}
  = \begin{bmatrix}w_{i1}\\ \vdots\\ w_{iV}\end{bmatrix} = (W_{i\cdot})^\top,
\]
the transpose of row $i$. Row $i$ of $W$ is the vector of logits used after the word $i$, and nothing
else: each row is a free parameter vector for one context. Implementations never form
the one-hot vector; they index the row (\texttt{W[i]}), which costs $V$ operations instead of $V^2$.

\begin{figure}[ht]
\centering
\begin{tikzpicture}[
  neu/.style={draw, circle, minimum size=7mm, inner sep=0pt, font=\scriptsize},
  ar/.style={-{Stealth[length=1.5mm]}, gray!70},
  arb/.style={-{Stealth[length=1.8mm]}, thick, bookaccent}]
  \foreach \k/\lab/\v in {1/{\langle s\rangle}/0,2/\text{the}/1,3/\text{a}/0,5/{\langle/s\rangle}/0} {
    \node[neu] (in\k) at (0,-0.95*\k) {$\v$};
    \node[font=\scriptsize, anchor=east] at (-0.5,-0.95*\k) {$\lab$};
  }
  \node[neu, fill=bookaccent!25] at (in2) {$1$};
  \node at (0,-0.95*4) {$\vdots$};
  \foreach \k/\lab in {1/{\langle s\rangle},2/\text{the},3/\text{a},5/{\langle/s\rangle}} {
    \node[neu, fill=bookexample!15] (out\k) at (3.6,-0.95*\k) {$z_{\k}$};
    \node[font=\scriptsize, anchor=west] at (5.4,-0.95*\k) {$p(\lab\mid\text{the})$};
    \draw[-{Stealth[length=1.5mm]}] (out\k) -- (5.3,-0.95*\k);
  }
  \node at (3.6,-0.95*4) {$\vdots$};
  \foreach \a in {1,3,5} \foreach \b in {1,2,3,5} \draw[ar] (in\a) -- (out\b);
  \foreach \b in {1,2,3,5} \draw[arb] (in2) -- (out\b);
  \node[font=\scriptsize, text=bookaccent] at (1.8,0.15) {row ``the'' of $W$};
  \node[draw, rounded corners=2pt, font=\scriptsize, fill=gray!8] at (4.6,-0.95*3) {softmax};
\end{tikzpicture}
\caption{Softmax regression for the bigram model. With the one-hot input of ``the'', only the
weights leaving the active unit (blue), that is row ``the'' of $W$, contribute to the logits.}
\label{bigram-models:fig:network}
\end{figure}

\paragraph{The bias is redundant.} Since $\sum_ix_i = 1$ for a one-hot vector, $W^\top x + b = (W + \mathbf 1b^\top)^\top x$, where
$\mathbf 1b^\top$ adds $b^\top$ to every row of $W$. Any model with a bias equals a model without bias, so the bias does not
enlarge the set of conditional tables the model can represent. It is harmless and can
help optimization (it learns the part of the logits shared by all contexts, the unigram
frequencies), but it is not needed; many implementations drop it. The softmax has a second
redundancy: adding the same constant to all the logits of a row does not change $p$.

\section{The loss and its gradient}\label{bigram-models:sec:gradient}

\paragraph{The loss.} For one bigram $(i, j)$ the loss is the negative log-likelihood of the observed
next word,
\[
  \ell = -\log p_j = -\sum_{k=1}^Vy_k\log p_k,\qquad y = e_j,
\]
the cross-entropy between the one-hot target $y$ and the prediction $p$.

\begin{pitfall}[The one-hot target is not an approximation]
It is often said that the ``true'' distribution of the word after ``the'' is not one-hot,
and that the one-hot target $y = e_j$ is an approximation made for convenience. It is not. The
target of one example is one-hot because one word was observed. Summed over the corpus,
the losses give \emph{exactly} the cross-entropy with the empirical distribution of the next word,
as the next proposition shows. Training on one-hot targets is maximum likelihood, and
maximum likelihood fits the full conditional distribution. Label smoothing, which replaces
$e_j$ by a mixture with the uniform distribution, is a regularizer, not a correction of an
approximation.
\end{pitfall}

\begin{proposition}[Empirical cross-entropy]\label{bigram-models:prop:empirical}
Let $\hat\pi_i = C_i/N$ and $\hat p(j\mid i) = C_{ij}/C_i$. For any model $p_\theta$,
\[
  \frac1N\sum_{n=1}^N-\log p_\theta(j_n\mid i_n) = \sum_{i:\,C_i>0}\hat\pi_i\,H\bigl(\hat p(\cdot\mid i),\,p_\theta(\cdot\mid i)\bigr),\qquad
  H(q, p) = -\sum_jq_j\log p_j .
\]
\end{proposition}

\begin{proof}
Group the $N$ terms by bigram: $(i, j)$ appears $C_{ij}$ times, so the left side is
\[
  -\frac1N\sum_{i,j}C_{ij}\log p_\theta(j\mid i) = -\sum_i\frac{C_i}N\sum_j\frac{C_{ij}}{C_i}\log p_\theta(j\mid i) . \qedhere
\]
\end{proof}

\paragraph{Gradient with respect to the logits.} With $\ell = -z_j + \log\sum_ke^{z_k}$,
\[
  \frac{\partial\ell}{\partial z_k} = -[k = j] + \frac{e^{z_k}}{\sum_{k'}e^{z_{k'}}} = p_k - y_k,\qquad\text{that is}\qquad\nabla_z\ell = p - y
\]
(\cref{backpropagation:sec:output}): predicted minus observed. Writing $\ell$ as $-\log p_j$ and differentiating
the softmax gives the same result through $\partial p_j/\partial z_k = p_j([k = j] - p_k)$, and the factor $p_j$ cancels against the
derivative $-1/p_j$ of the logarithm.

\paragraph{Gradient with respect to $W$.} The logit $z_k = \sum_iw_{ik}x_i + b_k$ depends on the column $k$ of $W$, and
$\partial z_k/\partial w_{ik'} = x_i[k = k']$. By the chain rule,
\[
  \frac{\partial\ell}{\partial w_{ik}} = \sum_{k'}\frac{\partial\ell}{\partial z_{k'}}\frac{\partial z_{k'}}{\partial w_{ik}} = x_i\,(p_k - y_k),
\]
so that, in matrix form,
\[
  \nabla_W\ell = x\,(p - y)^\top = \begin{bmatrix}x_1\\ \vdots\\ x_V\end{bmatrix}\begin{bmatrix}p_1 - y_1 & \cdots & p_V - y_V\end{bmatrix},\qquad
  \nabla_b\ell = p - y .
\]
This is the outer product of \cref{backpropagation:sec:affine}, transposed because the layer uses $W^\top$.
For $x = e_i$ only row $i$ is non-zero: one example updates only the logits of its own context,
\[
  W_{i\cdot}\leftarrow W_{i\cdot} - \eta\,(p - e_j)^\top .
\]
It raises the logit of the observed word by $\eta(1 - p_j)$ and lowers every other logit $k$ by $\eta p_k$.

\begin{example}[One gradient step]\label{bigram-models:ex:step}
In the corpus of \cref{bigram-models:ex:corpus}, start from $W = 0$, $b = 0$: every prediction is uniform,
$p_k = 1/8$, and the loss of the bigram (the, cat) is $\log8 = 2.0794$. The gradient row is
\[
  p - e_{\text{cat}} = (0.125,\ 0.125,\ 0.125,\ -0.875,\ 0.125,\ 0.125,\ 0.125,\ 0.125),
\]
in the order $\langle s\rangle$, the, a, cat, dog, sat, ran, $\langle/s\rangle$. One step with $\eta = 1$ sets the row ``the'' to $(-0.125,\dots,0.875,\dots,-0.125)$
and gives $p(\text{cat}\mid\text{the}) = e^{0.875}/(e^{0.875} + 7e^{-0.125}) = 0.2797$, up from $0.125$.
\end{example}

\section{What training converges to}\label{bigram-models:sec:convergence}

Summing the gradients of all the bigrams that start with $i$ gives the full-batch gradient of the
average loss with respect to row $i$:
\[
  \frac1N\sum_{n:\,i_n = i}(p_\theta(\cdot\mid i) - e_{j_n}) = \frac{C_i}N\bigl(p_\theta(\cdot\mid i) - \hat p(\cdot\mid i)\bigr).
\]
It vanishes exactly when the model reproduces the counts. Together with
\cref{bigram-models:prop:empirical}, this gives the following.

\begin{proposition}[Softmax regression learns the counts]\label{bigram-models:prop:same}
For every context $i$ with $C_i > 0$ whose counts are all positive ($C_{ij} > 0$ for all $j$), the loss is
minimized exactly when $\softmax(W_{i\cdot}^\top + b) = \hat p(\cdot\mid i)$, that is, when the logits are $\log\hat p(\cdot\mid i)$ up to an
additive constant. If some $C_{ij} = 0$, the infimum of the loss is the loss of the count model,
but it is not attained: driving the probability of the unseen bigram to $0$ requires $z_j\to-\infty$.
\end{proposition}

\begin{proof}
By \cref{bigram-models:prop:empirical} and Gibbs' inequality, the loss is at least that of the count
model, with equality iff $p_\theta(\cdot\mid i) = \hat p(\cdot\mid i)$ for every observed context. The softmax reaches a distribution $q$
iff $q$ has no zero entry, with logits $\log q$ plus a constant. If $\hat p(j\mid i) = 0$, the distributions $\softmax(z)$ with $z_j$
decreasing to $-\infty$ and the other logits fixed at $\log\hat p(\cdot\mid i)$ approach $\hat p(\cdot\mid i)$, so the loss approaches its lower
bound without reaching it.
\end{proof}

So the one-layer softmax model with a one-hot input is not a different model from the
count model: it is the count model, parameterized by logits and fitted by gradient descent.
Its interest lies elsewhere: it is the starting point for models that \emph{share} parameters
between contexts (\cref{ch:bigram-neural}), which counting cannot do.

\paragraph{Unseen bigrams and regularization.} The count model gives probability $0$ to every
bigram absent from the training data, so a test sentence with one new bigram has
probability $0$. In the corpus above, \emph{a dog ran} contains (a, dog) and (dog, ran), never seen. Two
remedies:
\begin{itemize}[nosep]
  \item \emph{smoothing} the counts, for instance add-one (Laplace) smoothing,
        $\hat p(j\mid i) = (C_{ij} + 1)/(C_i + V)$, or more refined methods (Kneser--Ney);
  \item \emph{regularizing} the softmax model: minimizing $\frac1N\sum_n\ell_n + \frac\lambda2\norm W^2$ keeps the logits finite, and
        the unseen words of each row keep a positive probability.
\end{itemize}

\paragraph{Perplexity.} A language model is evaluated by the average negative log-likelihood
per word on held-out text, or by its exponential, the \emph{perplexity}:
\[
  \mathrm{PPL} = \exp\Bigl(-\frac1N\sum_{n=1}^N\log p_\theta(j_n\mid i_n)\Bigr) .
\]
A model that is uniform over $K$ words has perplexity $K$: perplexity is the size of a uniform
choice that would be as uncertain as the model. A single zero probability makes it infinite.

\begin{example}[Fitting, overfitting, smoothing]\label{bigram-models:ex:fit}
Train the softmax model on the corpus of \cref{bigram-models:ex:corpus} by full-batch gradient descent
($\eta = 1$, $5\,000$ steps, $W$ initialized at $0$, no bias).
\begin{itemize}[nosep]
  \item With $\lambda = 0$, $p(\cdot\mid\text{the})$ is $0.6659$ for cat, $0.3325$ for dog and $0.0003$ for each of the six other
        words: almost the counts $\frac23$, $\frac13$, $0$. The training perplexity is $1.464$, close to $1.461$ for the
        count model, and the largest weight is $7.8$, still growing.
  \item With $\lambda = 0.01$, $p(\cdot\mid\text{the})$ is $0.549$, $0.256$ and $0.033$ for each other word, and the training
        perplexity rises to $1.860$.
\end{itemize}
On the held-out sentence \emph{a dog ran}, the count model has infinite perplexity; the softmax model
has perplexity $70.2$ with $\lambda = 0$ (and it keeps growing with training) and $7.45$ with $\lambda = 0.01$; add-one
smoothing gives $6.84$. Regularization trades training fit for robustness to unseen events,
exactly as in \cref{ch:erm}.
\end{example}

\begin{pitfall}[Tokenization]
The vocabulary is defined by the tokenizer, and the model knows nothing else. If the text is
split on spaces without lowercasing, ``To'' and ``to'' are two different words, with separate rows
and columns, and their statistics are split; punctuation glued to words (``be,'') creates further
spurious words. Decide on lowercasing, punctuation and sentence boundaries before
counting, and apply the same rules to the test data.
\end{pitfall}

\section{Implementation}\label{bigram-models:sec:code}

The training loop on single bigrams, with the gradient of \cref{bigram-models:sec:gradient}, in
NumPy. The softmax subtracts the maximum logit before exponentiating, which leaves the result
unchanged (the shift invariance above) and prevents overflow.
\begin{lstlisting}[language=Python]
import numpy as np

def softmax(z):
    e = np.exp(z - z.max())
    return e / e.sum()

def sgd_epoch(W, bigrams, lr):              # bigrams: list of index pairs (i, j)
    loss = 0.0
    for i, j in bigrams:
        p = softmax(W[i])                   # W.T @ one_hot(i) is the row W[i]
        loss -= np.log(p[j])
        g = p.copy(); g[j] -= 1.0           # dloss/dz = p - e_j
        W[i] -= lr * g                      # only row i receives a gradient
    return loss / len(bigrams)
\end{lstlisting}
In PyTorch the same model is \texttt{nn.Embedding(V, V)} followed by
\texttt{F.cross\_entropy}, which combines the log-softmax and the negative log-likelihood in a
numerically stable way and takes the target as a class index, not as a one-hot vector.

\begin{chaptersummary}
\begin{itemize}
  \item A language model is an autoregressive model of the next word; the bigram model
        conditions on the previous word only, and its maximum-likelihood estimate is the table
        of conditional frequencies $C_{ij}/C_i$.
  \item Softmax regression on a one-hot input computes $z = W^\top e_i + b$, the transpose of row $i$ of $W$: one free
        logit vector per context. The bias is redundant.
  \item $\nabla_z\ell = p - y$ and $\nabla_W\ell = x(p - y)^\top$; one example updates only the row of its context.
  \item Summed over the corpus, the one-hot cross-entropies equal the cross-entropy with the
        empirical conditional distributions: the one-hot target is exact, not an approximation.
  \item Training converges to the counts; unseen bigrams push logits to $-\infty$, which smoothing or
        regularization prevents. Perplexity is the exponential of the average negative
        log-likelihood.
\end{itemize}
\end{chaptersummary}

\section*{Exercises}
\addcontentsline{toc}{section}{Exercises}

\begin{exercise}\label{bigram-models:exr:1}
In \cref{bigram-models:ex:corpus}, compute the probability of the sentence \emph{the cat ran} under the count
model, and the training perplexity $1.461$.
\solutionlink{bigram-models:exr:1}
\end{exercise}

\begin{exercise}\label{bigram-models:exr:2}
Show that replacing $(W, b)$ by $(W + \mathbf 1c^\top,\ b - c)$ does not change any prediction, and that neither
does adding a constant to every entry of one row of $W$. How many free conditional probabilities
does the bigram model have, compared with the $V^2 + V$ parameters of $(W, b)$?
\solutionlink{bigram-models:exr:2}
\end{exercise}

\begin{exercise}\label{bigram-models:exr:3}
Compute $\nabla_W\ell$ and $\nabla_b\ell$ for $V = 3$, input $x = e_1$, prediction $p = (0.1, 0.7, 0.2)$ and observed next word $2$.
Which entries of $W$ change in a gradient step?
\solutionlink{bigram-models:exr:3}
\end{exercise}

\begin{exercise}\label{bigram-models:exr:4}
With label smoothing, the target is $y = (1 - \varepsilon)e_j + \frac\varepsilon V\mathbf 1$. Give $\nabla_z\ell$, and the logits that minimize
the loss for a context that is always followed by the same word $j$.
\solutionlink{bigram-models:exr:4}
\end{exercise}

\begin{exercise}\label{bigram-models:exr:5}
A model is uniform over the $V$ words. Show that its perplexity is $V$ on any text. Show that the
perplexity of any model on a text is at least $1$, with equality iff it gives probability $1$ to every
observed word.
\solutionlink{bigram-models:exr:5}
\end{exercise}
\end{notechapter}
